% ================================================================
% SECTION 2 — LITERATURE REVIEW AND RESEARCH GAP
% Target: ~7–8 pages
% ================================================================

\section{Literature Review and Research Gap}
\label{sec:literature}

\subsection{Causal AI for Time-Series Data}
\label{sec:causal_ts}

The term \emph{causal AI} refers to the integration of causal
reasoning into data-driven modelling --- the pursuit of methods that
can answer not just ``what is correlated with the outcome?'' but
``what would happen to the outcome if we intervened on a
variable?''~\parencite{pearl2009, scholkopf2021}.
In the time-series setting, this distinction is especially
consequential because temporal ordering creates a natural temptation
to treat precedence as causation~\parencite{granger1969}.

At the most basic level, \emph{Granger causality}~\parencite{granger1969}
formalises the intuition that if past values of variable $X$ help
predict variable $Y$ even after conditioning on $Y$'s own past, then
$X$ ``Granger-causes'' $Y$.
Granger causality is a statistical property of predictive precedence,
not a structural causal claim~\parencite{granger1969}.
A variable can Granger-cause another purely because both are driven by
a common latent factor, or because the causal direction is actually
reversed but operates at a different lag~\parencite{pearl2009}.
The existence of a Granger-causal relationship is therefore neither
necessary nor sufficient for the existence of a structural causal
relationship in the interventionist sense~\parencite{pearl2009}.

\emph{Structural causal interpretation} goes further.
It requires specifying a causal graph --- a directed acyclic graph
(DAG) or, in the temporal case, a time-unrolled directed graph ---
that encodes the researcher's assumptions about which variables affect
which others, and with what lag~\parencite{pearl2009}.
Under this framework, the causal effect of a treatment $T$ on an
outcome $Y$ is defined as the change in $Y$ under an intervention
that sets $T$ to a specific value, holding all other structural
equations constant~\parencite{pearl2009}.
This is the \emph{do-calculus} of \textcite{pearl2009}, and it
is fundamentally different from conditioning on observed values of $T$.
The distinction has been formalised in the potential outcomes
framework~\parencite{rubin1974} and is central to the
identification strategy used in this thesis.

Time ordering helps causal identification because a cause must
precede its effect, which rules out reverse causation within a given
lag window~\parencite{granger1969, pearl2009}.
However, time ordering does not automatically solve the confounding
problem~\parencite{angrist2009}.
A time-series confounder --- a variable that drives both $T_t$ and
$Y_{t+k}$ --- induces a spurious association between treatment and
outcome regardless of the temporal structure of the data.
In the freight rail context, load-shedding hours are a clear
candidate: they simultaneously reduce locomotive availability
(affecting rail supply) and depress industrial production (affecting
freight demand)~\parencite{havenga2021, transnet2023ar}.
Any analysis that does not adequately account for such confounders
will attribute their effects to whichever treatment variable happens
to be correlated with them~\parencite{pearl2009}.

For this reason, the most credible causal analyses of time-series
interventions do not rely on prediction models or Granger tests
alone~\parencite{abadie2010, benmichael2021}.
They instead use counterfactual methods that construct an explicit
baseline --- what the outcome would have been absent the
intervention --- and identify the treatment effect as the deviation
between the observed outcome and that
baseline~\parencite{rubin1974, abadie2003}.
The Synthetic Control Method, reviewed in the next section, is the
leading such method for settings with a single treated unit observed
over time~\parencite{abadie2010}.

\subsection{The Synthetic Control Family}
\label{sec:scm_family}

\subsubsection{Standard Synthetic Control Method}
\label{sec:base_scm}

The Synthetic Control Method (SCM) was introduced by
\textcite{abadie2003} to evaluate the economic
impact of terrorism in the Basque Country, and formalised by
\textcite{abadie2010} in the context
of California's tobacco control programme.
The core idea is to construct a weighted average of untreated
``donor'' units whose pre-treatment outcome trajectory closely
matches that of the treated unit~\parencite{abadie2010}.
This weighted average --- the \emph{synthetic control} --- then
serves as the counterfactual for the post-treatment period.
The method has since been applied across a wide range of settings
including natural disasters~\parencite{cavallo2013}, trade
policy~\parencite{abadie2015}, and public health
interventions~\parencite{abadie2010}.

Formally, let $Y_{it}$ denote the outcome for unit $i$ at time $t$,
with unit $i = 1$ being the treated unit and units
$i = 2, \ldots, J+1$ forming the donor pool.
Let $T_0$ denote the treatment date, so the pre-period is
$t = 1, \ldots, T_0$ and the post-period is
$t = T_0+1, \ldots, T$.
The SCM solves for a weight vector
$\mathbf{w} = (w_2, \ldots, w_{J+1})^\top$ with $w_j \geq 0$ and
$\sum_j w_j = 1$ by minimising the pre-period outcome discrepancy:
\begin{equation}
  \hat{\mathbf{w}} = \argmin_{\mathbf{w} \in \Delta^J}
  \sum_{t=1}^{T_0} \!\left(Y_{1t} - \sum_{j=2}^{J+1} w_j Y_{jt}\right)^{\!2}
  \label{eq:scm}
\end{equation}
where $\Delta^J$ is the unit simplex.
The causal effect estimate for each post-period $t > T_0$ is then:
\begin{equation}
  \hat{\tau}_t = Y_{1t} - \sum_{j=2}^{J+1} \hat{w}_j Y_{jt}
  \label{eq:gap}
\end{equation}

The identifying assumption is that, in the absence of treatment, the
outcome for the treated unit would have evolved in the same way as
the synthetic control --- a form of the parallel-trends assumption
extended to allow for a flexible, data-driven weighting of the
comparison group~\parencite{abadie2010}.
\textcite{abadie2010} show that this assumption is supported when
pre-period fit is tight and the donor pool is broad enough to span
the treated unit's trajectory.

The standard SCM has two well-known limitations~\parencite{benmichael2021}.
First, the convex hull constraint ($\mathbf{w} \in \Delta^J$) can
produce poor fits when the treated unit lies outside the support of
the donor pool.
Second, with a small donor pool (in the extreme case, a single donor),
the weight vector is trivially determined and the synthetic control
is nothing more than a rescaled version of the single donor ---
a counterfactual that is unlikely to be credible~\parencite{abadie2003}.
Both limitations are directly relevant to this thesis and motivate the
extended estimators reviewed below.

\subsubsection{Augmented Synthetic Control Method (ASCM)}
\label{sec:ascm}

\textcite{benmichael2021} introduced
the Augmented Synthetic Control Method to address the bias that arises
when the standard SCM fit is imperfect.
ASCM adds a ridge-regression bias-correction term on top of the
standard SCM counterfactual.
Let $\hat{Y}_{1t}^{\text{SCM}} = \sum_j \hat{w}_j Y_{jt}$ be the
standard SCM counterfactual and
$\hat{\varepsilon}_t = Y_{1t} - \hat{Y}_{1t}^{\text{SCM}}$
be the pre-period residuals.
ASCM learns a ridge model $\hat{m}_\alpha(\cdot)$ from the pre-period
residuals and applies it to the post-period donor features:
\begin{equation}
  \hat{Y}_{1t}^{\text{ASCM}} = \hat{Y}_{1t}^{\text{SCM}}
    + \hat{m}_\alpha\!\left(\mathbf{X}_t\right)
  \label{eq:ascm}
\end{equation}
where $\mathbf{X}_t$ contains donor values and time features, and
$\alpha$ is selected by blocked cross-validation that respects the
temporal ordering of observations~\parencite{benmichael2021}.
When the pre-period fit is already excellent (as in the five-donor
analysis), the ridge correction is small and ASCM converges toward
the standard SCM estimate.
When the pre-period fit is poor (as in the single-donor case), the
correction can be large but potentially unstable~\parencite{benmichael2021}.

\subsubsection{Elastic Net SCM (EN-SCM)}
\label{sec:enscm}

Rather than correcting for imperfect fit after the fact, EN-SCM
replaces the hard simplex constraint in Equation~\ref{eq:scm} with
a penalised objective that combines $L_1$ (Lasso) and $L_2$ (Ridge)
penalties on the donor weights~\parencite{hollingsworth2020}.
The elastic net penalty was originally proposed by
\textcite{zouelastic2005} as a compromise between the sparsity of Lasso
and the stability of Ridge regression, and its application to
synthetic control settings was developed by
\textcite{hollingsworth2020}:
\begin{equation}
  \hat{\mathbf{w}}^{\text{EN}} = \argmin_{\mathbf{w} \geq 0}
  \sum_{t=1}^{T_0}\!\left(Y_{1t} - \sum_j w_j Y_{jt}\right)^{\!2}
  + \alpha\!\left[\frac{1-\rho}{2}\|\mathbf{w}\|_2^2
    + \rho\|\mathbf{w}\|_1\right]
  \label{eq:enscm}
\end{equation}
where $\alpha > 0$ controls overall regularisation strength and
$\rho \in [0,1]$ controls the $L_1/L_2$ balance.
The weights are non-negative but need not sum to one.
When $\rho = 0$ the penalty is pure ridge; when $\rho = 1$ it is
pure Lasso, which induces sparsity~\parencite{zouelastic2005}.
Both $\alpha$ and $\rho$ are selected jointly by blocked
cross-validation.
EN-SCM is particularly useful when donors are correlated and a sparse,
stable weight vector is preferable to a uniform
one~\parencite{hollingsworth2020}.

\subsubsection{Matrix Completion SCM (MC-SCM)}
\label{sec:mcscm}

\textcite{athey2021} reframe counterfactual estimation as
a matrix completion problem.
The full outcome panel $\mathbf{Y} \in \mathbb{R}^{T \times N}$ is
treated as a low-rank matrix, and the post-treatment entries of the
treated unit are treated as missing cells to be imputed~\parencite{athey2021}.
Recovery proceeds via Singular Value Thresholding (SVT), a technique
developed by \textcite{cai2010svt} for low-rank matrix recovery from
incomplete observations:
iteratively decompose the partially-observed matrix, shrink singular
values by a nuclear norm penalty $\lambda$, and re-pin the observed
entries to their true values.
The algorithm converges to a low-rank completion that minimises:
\begin{equation}
  \min_{\mathbf{M}} \;
  \frac{1}{|\Omega|}\sum_{(i,t)\in\Omega}(Y_{it} - M_{it})^2
  + \lambda \|\mathbf{M}\|_*
  \label{eq:mcscm}
\end{equation}
where $\Omega$ is the set of observed entries and $\|\cdot\|_*$
is the nuclear norm (sum of singular values)~\parencite{athey2021}.
MC-SCM does not use explicit donor weights; instead it exploits the
latent factor structure shared across all corridors.
With six corridors over 120 months, the panel has enough
cross-sectional and temporal variation for the low-rank assumption
to be meaningful, as confirmed by an effective rank of four in the
five-donor analysis.

\subsubsection{Synthetic Difference-in-Differences (SDID)}
\label{sec:sdid}

\textcite{arkhangelsky2021} introduce SDID as a
method that simultaneously solves for two sets of weights:
\emph{unit weights} $\boldsymbol{\lambda}$ (identical to standard
SCM weights) that balance donor trajectories to the treated unit
in the pre-period, and \emph{time weights} $\boldsymbol{\omega}$
that re-weight pre-period time points to resemble the post-period
donor distribution.
SDID can be seen as a synthesis of the Difference-in-Differences
estimator~\parencite{card1994} and the standard SCM, combining the
double-differencing logic of DiD with the data-driven weighting of
SCM~\parencite{arkhangelsky2021}.
The SDID estimator is:
\begin{equation}
  \hat{\tau}^{\text{SDID}} =
    \left(\bar{Y}_{1,\text{post}} - \sum_t \omega_t Y_{1t}\right)
    - \sum_j \lambda_j \!\left(\bar{Y}_{j,\text{post}}
      - \sum_t \omega_t Y_{jt}\right)
  \label{eq:sdid}
\end{equation}
The time weights remove pre-period time shocks --- such as the
COVID-19 disruption of 2020--2021 --- that affected all corridors
uniformly and do not persist into the post-period~\parencite{arkhangelsky2021}.
This makes SDID more robust to common temporal confounders than
standard SCM or DiD alone~\parencite{arkhangelsky2021}.
Standard errors are computed via the jackknife, dropping one
pre-treatment year at a time~\parencite{arkhangelsky2021}.

\subsection{Causal Effect Estimation in Transport and Infrastructure}
\label{sec:transport_causal}

The application of formal causal inference methods to transport and
infrastructure economics is growing but remains limited in the
African context~\parencite{jedwab2016}.
Studies of natural disasters and infrastructure disruptions in other
settings provide useful methodological benchmarks.

\textcite{cavallo2013} used SCM to estimate the
macroeconomic impact of large natural disasters, finding that
catastrophic events can have persistent negative effects on GDP that
diverge substantially from naive pre-post comparisons.
Their finding that a single poorly-chosen donor can produce wildly
different estimates from a well-constructed pool is directly relevant
to the present study.
Similarly, \textcite{barone2014} applied difference-in-differences
to estimate the effect of earthquake damage on local economic activity
in Italy, finding persistent effects lasting five or more years ---
consistent with the hysteresis hypothesis tested in this thesis.

\textcite{abadie2003} applied SCM to estimate the economic cost of
terrorism in the Basque Country, finding a GDP loss of approximately
10 percentage points relative to the synthetic control --- a result
that survived in-space placebo tests using other Spanish regions as
pseudo-treated units.
Their paper established the placebo-based inference framework that
all subsequent SCM applications, including this thesis, follow.

\textcite{abadie2010} applied SCM to California's 1988 tobacco
control programme using other US states as donors.
This paper introduced the formal leave-one-out and in-space placebo
procedures and demonstrated that SCM estimates are most credible when
the pre-period fit is tight and the donor pool is
broad~\parencite{abadie2010}.
Both conditions are satisfied in the five-donor analysis of this
thesis but violated in the single-donor version.

In the transport infrastructure context more broadly,
\textcite{jedwab2016} examine the long-run causal effects of colonial
railway construction on African urbanisation, using geographic
discontinuities as instruments.
Their work illustrates both the importance of causal identification
in transport research and the methodological challenges of applying
it when treatment assignment is non-random~\parencite{jedwab2016}.
\textcite{jedwab2016} further demonstrate that transport
infrastructure effects in Africa are highly persistent and
heterogeneous across regions, with long-run economic impacts
far exceeding short-run estimates.

In the South African freight rail context, \textcite{stander2005}
provide a comprehensive descriptive analysis of the road--rail modal
split and the structural drivers of rail's declining market share
since deregulation.
They document how rail's share of outsourced freight fell from
approximately 53\% in 1957 to around 20--25\% by the late 1980s
and continued to erode thereafter~\parencite{stander2005}.
Crucially, their analysis is entirely descriptive and
correlation-based: they identify co-movement between freight volumes
and macroeconomic conditions but make no causal claims.

\textcite{havenga2021} develop a road-to-rail strategy for South
Africa using the national Freight Demand Model (FDMTM), establishing
that rail's tonne-km market share dropped from 23\% in 2011 to 16\%
in 2019 and that the system was approaching a critical threshold of
infrastructure deterioration~\parencite{havenga2021}.
They frame the problem in terms of macrologistics policy and
investment strategy but do not employ causal identification methods.
Like \textcite{stander2005}, their analysis cannot attribute specific
volume changes to specific causes, and the report explicitly calls
for more rigorous data-driven policy analysis~\parencite{havenga2021}.

\subsection{Domain Literature Benchmark Table}
\label{sec:benchmark}

Table~\ref{tab:benchmark} positions this thesis against a selection
of representative studies spanning both the SCM methodology literature
and the South African freight domain.
The table highlights the dimensions along which prior work falls short
and where this thesis contributes.

\begin{table}[H]
  \centering
  \caption{Benchmark of representative studies. ``Causal discovery''
    refers to explicit identification of a causal graph or treatment
    effect; ``robustness'' refers to placebo, leave-one-out, or
    sensitivity tests. SA = South Africa.}
  \label{tab:benchmark}
  \small
  \renewcommand{\arraystretch}{1.35}
  \begin{tabularx}{\textwidth}{%
      >{\raggedright\arraybackslash}p{3.2cm}   % Study
      >{\raggedright\arraybackslash}p{2.0cm}   % Domain
      >{\raggedright\arraybackslash}p{1.5cm}   % Target
      >{\raggedright\arraybackslash}p{1.3cm}   % Granularity
      >{\raggedright\arraybackslash}p{1.5cm}   % Method
      >{\centering\arraybackslash}p{1.1cm}     % Causal?
      >{\centering\arraybackslash}p{1.3cm}     % Robust?
      >{\raggedright\arraybackslash}p{2.5cm}   % Limitation
    }
    \toprule
    \textbf{Study} &
    \textbf{Domain} &
    \textbf{Target} &
    \textbf{Granularity} &
    \textbf{Method} &
    \textbf{Causal?} &
    \textbf{Robust?} &
    \textbf{Main limitation} \\
    \midrule
    \textcite{abadie2003} &
    Regional economics &
    GDP per capita &
    Annual &
    SCM &
    Yes &
    Yes &
    Single treated region; cross-sectional donors only \\

    \textcite{abadie2010} &
    Public health policy &
    Tobacco sales &
    Annual &
    SCM &
    Yes &
    Yes &
    Cross-sectional donors; no temporal weighting \\

    \textcite{abadie2015} &
    Comparative politics &
    GDP per capita &
    Annual &
    SCM &
    Yes &
    Yes &
    Cross-sectional donors; limited to annual data \\

    \textcite{arkhangelsky2021} &
    Methodology &
    Multiple &
    Varies &
    SDID &
    Yes &
    Yes &
    Requires balanced panel; many donors preferred \\

    \textcite{benmichael2021} &
    Methodology &
    Multiple &
    Varies &
    ASCM &
    Yes &
    Yes &
    Ridge correction can be unstable with few donors \\

    \textcite{athey2021} &
    Methodology &
    Multiple &
    Varies &
    MC-SCM &
    Yes &
    Yes &
    Requires sufficient cross-sectional variation \\

    \textcite{cavallo2013} &
    Disaster economics &
    GDP &
    Annual &
    SCM &
    Yes &
    Yes &
    Annual granularity; donor pool sensitivity not explored \\

    \textcite{jedwab2016} &
    Transport (Africa) &
    Urbanisation &
    Decadal &
    IV/RDD &
    Yes &
    Partial &
    Long-run only; no short-run shock analysis \\

    \textcite{stander2005} &
    SA freight &
    Modal share &
    Multi-year &
    Descriptive &
    No &
    No &
    Purely descriptive; no counterfactual; no treatment effect \\

    \textcite{havenga2021} &
    SA freight &
    Tonne-km share &
    Annual &
    Demand model &
    No &
    No &
    Demand modelling only; no causal identification \\

    \midrule
    \textbf{This study} &
    SA rail freight &
    Rail volume (MT) &
    Monthly &
    SCM family (5 estimators) &
    \textbf{Yes} &
    \textbf{Yes} &
    Single treated corridor; observational data only \\
    \bottomrule
  \end{tabularx}
\end{table}

The table reveals three clear patterns.
First, the causal inference literature has produced a mature set of
SCM-family tools, each with well-understood strengths and
limitations~\parencite{abadie2010, benmichael2021, arkhangelsky2021},
but their application to freight and logistics is limited.
Second, the South African freight literature is rich in descriptive
and demand-modelling analysis~\parencite{stander2005, havenga2021}
but lacks any study that applies a counterfactual identification
strategy to a specific policy shock or natural disaster.
Third, no prior study combines monthly time-series granularity,
multiple SCM estimators, and systematic robustness testing in the
South African freight rail context.
This thesis fills all three gaps simultaneously.

\subsection{Research Gap and Thesis Positioning}
\label{sec:gap_statement}

The foregoing review reveals a structured gap at the intersection of
causal inference methodology and South African freight rail analysis.

Prior studies that address causal questions do so either at annual
granularity (limiting the detection of short-run dynamics), with a
cross-sectional donor pool (limiting temporal
generalisability)~\parencite{abadie2003, abadie2010}, or in contexts
unrelated to transport infrastructure~\parencite{cavallo2013}.
Prior studies that address the South African freight context are
descriptive or predictive but not causal: they cannot distinguish
the effect of a specific shock from the concurrent influence of
macroeconomic trends, operational deterioration, or policy
changes~\parencite{stander2005, havenga2021}.

More specifically, the gap has four dimensions.
First, no existing study applies a counterfactual framework to
estimate the causal effect of a specific infrastructure shock on
South African rail freight volumes~\parencite{havenga2021}.
Second, no study examines the sensitivity of SCM estimates to donor
pool size in the freight context, despite this being a well-known
source of estimator instability~\parencite{abadie2003, benmichael2021}.
Third, no study applies more than one SCM-family estimator to the
same freight dataset, making it impossible to assess the robustness
of any single estimate across alternative modelling
assumptions~\parencite{benmichael2021, athey2021}.
Fourth, no study uses monthly panel data on South African freight
corridors over a sufficiently long horizon to support the
pre-period balance requirements of modern SCM
variants~\parencite{arkhangelsky2021}.

This thesis addresses all four dimensions.
It constructs a monthly panel of six freight corridors spanning
120 months (2015--2024), applies five SCM-family estimators under
a common evaluation protocol, and explicitly demonstrates how the
number of donor corridors determines the coherence of the estimates.
The result is not merely a point estimate of the KZN flood effect,
but a methodological argument: that causal analysis of freight
infrastructure shocks requires both an adequate donor pool and
multi-estimator robustness assessment to be credible~\parencite{benmichael2021}.
